% =============================================================================
%  Notas del Profesor — Sesión 21: Archivos de texto
%  Programación Avanzada — Universidad Panamericana
%  Compilar: pdflatex sesion21_archivos_texto_profesor.tex
% =============================================================================
\documentclass[12pt, oneside]{article}
\usepackage[profesor]{progavanzada}
\usepackage{array}

\begin{document}

\portada{Sesión 21}{Archivos de texto}{2026}

\tableofcontents
\newpage

% =============================================================================
\section{Objetivos de la sesión}
% =============================================================================

Al finalizar la sesión el alumno será capaz de:

\begin{enumerate}
  \item Explicar la diferencia entre almacenamiento en memoria RAM y en
        archivo de disco, y cuándo conviene cada uno.
  \item Crear, leer y escribir archivos de texto en C++ usando
        \cpp{ifstream} y \cpp{ofstream}.
  \item Usar \cpp{getline} para leer líneas completas y \cpp{>>} para
        leer tokens numéricos de un archivo.
  \item Escribir un programa que genere un archivo de datos estructurado,
        y otro que lo consuma para responder consultas.
  \item Participar en una actividad colaborativa donde los datos de varios
        programas se integran en un solo archivo.
\end{enumerate}

\begin{notaprofesor}[Duración y distribución (100 min)]
  \begin{tabular}{rl}
    10 min & Gancho histórico + preguntas de arranque \\
    10 min & Símil del escritorio; tabla RAM vs.\ disco \\
    15 min & ¿Qué es un archivo?  Demo con editor hex \\
     5 min & Presentar \cpp{<fstream>} y las tres clases \\
    25 min & Demo en vivo: archivos1, archivos2, archivos3 \\
    35 min & Actividad: agenda de cumpleaños (Parte~A + Parte~B) \\
  \end{tabular}
\end{notaprofesor}

% =============================================================================
\section{Prerrequisitos}
% =============================================================================

\begin{itemize}
  \item \cpp{cin} y \cpp{cout} con \cpp{>>} y \cpp{<<} (sesiones 1--2).
  \item \cpp{string} y \cpp{getline} desde teclado (sesión 2).
  \item Tabla ASCII y representación de caracteres (sesión 3).
  \item Compilación básica con \texttt{g++} y ejecución desde terminal.
\end{itemize}

\begin{notaprofesor}
  Esta sesión no requiere apuntadores, structs ni arreglos avanzados.  La
  actividad final sí puede llegar a necesitar arreglos simples o vectores
  para almacenar los datos de la agenda, pero los alumnos tienen esa base
  de sesiones anteriores.
  \par\smallskip
  Señal de alerta: si algún alumno no recuerda la diferencia entre
  \cpp{cin >>} y \cpp{getline}, detenerse 3 minutos a recordarlo con un
  ejemplo en pantalla antes de abrir archivos.  El modelo mental de
  ``el buffer del teclado'' es el mismo que el de ``el buffer del archivo''.
\end{notaprofesor}

% =============================================================================
\section{Arranque y preguntas de sondeo}
% =============================================================================

\begin{notaprofesor}[Preguntas de arranque (10 min)]

  \textbf{1.\ ``¿Qué le pasa a una variable cuando el programa termina?''}\\
  Respuesta esperada: desaparece / se pierde.  Seguir con: ``¿y si quieres
  que ese valor esté disponible la próxima vez que ejecutes el programa?''
  Dejar que propongan soluciones (variable global, retornar el valor, etc.)
  antes de decir: ``la solución real se llama archivo.''

  \medskip
  \textbf{2.\ ``¿Cuántas veces abrieron hoy un archivo antes de llegar
  aquí?''}\\
  Respuesta esperada: muchas (WhatsApp, fotos, música, documentos).
  Seguir con: ``todos esos archivos son secuencias de bytes.  Hoy van a
  crear sus propios archivos desde C++.''

  \medskip
  \textbf{3.\ ``¿Por qué creen que la carpeta del explorador de archivos
  se llama \textit{carpeta}?''}\\
  Respuesta esperada: porque se parece a una carpeta real.  Usar esto
  para introducir la historia de Xerox PARC y el origen deliberado de la
  metáfora.  Es el gancho para la sección histórica.

\end{notaprofesor}

% =============================================================================
\section{Memoria vs.\ archivos}
% =============================================================================

\begin{notaprofesor}[Cómo presentar el símil (5 min)]
  El símil del escritorio funciona bien si se actúa físicamente: señalar
  el escritorio (o la pantalla), abrir el cajón del escritorio real si hay
  uno.  Preguntar: ``¿qué pasa si empujas todos los papeles del escritorio
  al piso?''  Eso es lo que hace el sistema operativo cuando el programa
  termina: limpia la memoria.  ``¿Y el cajón?''  Los cajones se quedan como
  están.
  \par\smallskip
  Después mostrar la tabla comparativa y pedir a un alumno que la lea en
  voz alta.  Preguntar: ``¿cuándo prefieren usar memoria?  ¿cuándo un
  archivo?''  Respuesta esperada: memoria para cálculos en curso, archivo
  para guardar resultados permanentes.
\end{notaprofesor}

% =============================================================================
\section{¿Qué es un archivo?  Demo con hex editor}
% =============================================================================

\begin{notaprofesor}[Guión de la demo (15 min)]
  Esta demo es uno de los momentos más memorables de la sesión.  El objetivo
  es que los alumnos \textit{vean} que un archivo es solo bytes, no magia.

  \begin{enumerate}
    \item Crear un archivo de texto vacío:
          \begin{itemize}
            \item Windows: clic derecho → Nuevo → Documento de texto.
            \item Linux: \texttt{touch prueba.txt} en terminal.
          \end{itemize}
          Mostrar el tamaño: \textbf{0 bytes}.  Preguntar: ``¿qué esperaban?''

    \item Abrir con el editor de texto, escribir \texttt{H} y guardar.
          Mostrar el tamaño: \textbf{1 byte}.  Preguntar: ``¿cuántos bytes
          esperaban?''

    \item Abrir \texttt{prueba.txt} en \texttt{hexed.it} (o cualquier hex
          editor en línea).  Mostrar el byte \texttt{0x48} y conectarlo con
          la tabla ASCII de sesiones anteriores.

    \item Escribir \texttt{Hola} completo.  En el hex editor: cuatro bytes,
          \texttt{48 6F 6C 61}.  Pedir a algún alumno que identifique cada
          carácter.

    \item Agregar un salto de línea con Enter y guardar.  En Linux aparece
          \texttt{0A}; en Windows aparece \texttt{0D 0A}.  Explicar que
          \texttt{0A} es \texttt{LF} (Line Feed) y \texttt{0D} es
          \texttt{CR} (Carriage Return) — herencia de las máquinas de escribir
          donde había que hacer dos acciones separadas para pasar de línea.

    \item Preguntar: ``¿qué diferencia hay entre este archivo y una imagen
          JPEG?''  Respuesta: ninguna en fondo — también son bytes.  Lo que
          cambia es cómo el programa los interpreta.
  \end{enumerate}

  \textbf{Por qué vale la pena hacer esto:} muchos alumnos creen que un
  archivo de texto es algo cualitativamente distinto a un archivo binario.
  Ver los bytes directamente rompe esa idea y prepara la sesión~22 sobre
  archivos binarios.
\end{notaprofesor}

\begin{notaprofesor}[FAQ: archivos y bytes]

  \textbf{¿Por qué el archivo vacío tiene 0 bytes y no algún encabezado?}\\
  En sistemas Unix y la mayoría de sistemas modernos, un archivo de texto
  plano (\texttt{.txt}) no tiene encabezado.  Empieza directamente con los
  datos.  Otros formatos como JPEG, PDF o DOCX sí tienen encabezados
  (\textit{magic bytes}) que identifican el tipo de archivo.  Los primeros
  2 bytes de cualquier JPEG son siempre \texttt{FF D8}.

  \medskip
  \textbf{¿\texttt{sizeof(char)} puede ser distinto de 1?}\\
  No.  El estándar C++ garantiza \cpp{sizeof(char) == 1} por definición.
  Otras preguntas sobre el tamaño de los tipos son válidas (\cpp{int} puede
  ser 2 o 4 bytes según la plataforma), pero \cpp{char} siempre es 1.

  \medskip
  \textbf{¿El problema CRLF vs LF les va a afectar en este curso?}\\
  Raramente, porque \cpp{<fstream>} en modo texto hace la conversión
  automáticamente al leer y al escribir.  El problema surge cuando se
  abre el archivo en modo binario o cuando se comparte entre sistemas
  sin conversión.  Para la actividad de hoy no es un problema.

\end{notaprofesor}

% =============================================================================
\section{La biblioteca \texttt{<fstream>}}
% =============================================================================

\begin{notaprofesor}[Presentación (5 min)]
  Al mostrar la tabla de las tres clases, enfatizar la simetría con
  \cpp{cin}/\cpp{cout}:
  \begin{itemize}
    \item \cpp{cin} es un \cpp{istream} conectado al teclado.
    \item \cpp{ifstream} es un \cpp{istream} conectado a un archivo.
    \item La diferencia es solo dónde vienen los datos; los operadores
          son idénticos.
  \end{itemize}
  Esto reduce la carga cognitiva: los alumnos ya saben usar \cpp{>>} y
  \cpp{<<}.  Solo aprenden a apuntar esos operadores a un archivo.
\end{notaprofesor}

% =============================================================================
\section{Demo en vivo: los tres ejemplos}
% =============================================================================

\subsection{\texttt{archivos1.cpp}: leer línea por línea}

\begin{notaprofesor}[Guión (8 min)]
  \begin{enumerate}
    \item Mostrar \texttt{canciones.txt} con \texttt{cat} (Linux) o
          abrirlo en el editor.  Confirmar que tiene 3 líneas.
    \item Abrir \texttt{archivos1.cpp}, leerlo en voz alta línea por línea
          y preguntar antes de ejecutar: ``¿qué esperan que imprima?''
    \item Compilar y ejecutar.  Confirmar la salida.
    \item Preguntar: ``¿qué pasaría si \texttt{canciones.txt} no existiera?''
          Cambiar el nombre del archivo a algo que no existe y mostrar
          que el programa termina sin imprimir nada ni dar error visible.
    \item Agregar la verificación \cpp{if (!ifile)} justo después de abrir
          el archivo y repetir.  Ahora sí se muestra el mensaje de error.
          Esta es la forma correcta de manejar archivos que pueden no existir.
    \item Pedir a un alumno que explique qué hace la condición del
          \cpp{while (getline(...))}.  Respuesta esperada: \cpp{getline}
          devuelve el stream; cuando llega al EOF el stream es falso.
  \end{enumerate}
\end{notaprofesor}

\subsection{\texttt{archivos2.cpp}: escribir en un archivo}

\begin{notaprofesor}[Guión (8 min)]
  \begin{enumerate}
    \item Antes de ejecutar, preguntar: ``¿dónde va a aparecer la salida?''
          Respuesta esperada: en \texttt{resultados.txt}, no en pantalla.
    \item Ejecutar.  Abrir \texttt{resultados.txt} con \texttt{cat} y
          mostrar el contenido.
    \item Preguntar: ``¿cuántos bytes tiene ese archivo?''
          \texttt{ls -lh resultados.txt} o abrir en el hex editor.
          El archivo tiene \texttt{30 10 15\textbackslash n} = 9 caracteres = 9 bytes
          (o 10 en Windows por el \texttt{\textbackslash r}).
    \item Ejecutar el programa \textit{otra vez} sin cambiar nada.
          Mostrar que \texttt{resultados.txt} sigue igual: \cpp{ofstream}
          sobreescribe.  Introducir \cpp{fstream::app} y demostrar que
          ahora agrega una segunda línea.
    \item Preguntar: ``¿qué pasa si no llamamos a \cpp{close()}?''
          En la mayoría de los casos el destructor del objeto lo cierra
          automáticamente al salir del \cpp{main}, pero es mala práctica
          depender de eso, especialmente en programas largos donde el
          archivo se puede necesitar antes de que termine el programa.
  \end{enumerate}
\end{notaprofesor}

\subsection{\texttt{archivos3.cpp}: leer valores numéricos}

\begin{notaprofesor}[Guión (9 min)]
  \begin{enumerate}
    \item Mostrar que este programa lee lo que \textbf{archivos2} escribió.
          Enfatizar: comunicación entre programas a través del archivo.
    \item Antes de ejecutar: ``¿cuántos \cpp{>>} necesitamos para leer
          los tres números?  ¿Pueden ir en la misma línea?''
          Mostrar ambas formas (separada y encadenada) y confirmar que
          producen el mismo resultado.
    \item Ejecutar y verificar la salida.
    \item Pregunta de profundidad: ``Si el archivo tuviera más de 3 números,
          ¿qué pasaría?''  Respuesta: los \cpp{>>} leen solo los primeros 3;
          el resto queda sin leer en el stream.  ``¿Y si tuviera menos de
          3?''  Los \cpp{>>} fallan silenciosamente y \cpp{x}, \cpp{y} o
          \cpp{z} quedan con valores basura.  Para detectarlo se puede
          verificar el estado del stream: \cpp{if (!fresultados)} después
          de las lecturas.
    \item Conectar con la actividad: ``El programa que van a escribir en
          Parte~B lee datos estructurados de un archivo.  Ese es exactamente
          el patrón de archivos3.''
  \end{enumerate}
\end{notaprofesor}

\begin{notaprofesor}[FAQ: streams y operadores]

  \textbf{¿Por qué \cpp{getline} y no \cpp{>>} para leer canciones?}\\
  \cpp{>>} lee hasta el primer espacio o salto de línea.  ``Walking on
  Sunshine'' tiene espacios, así que \cpp{>>} leería solo ``Walking''.
  \cpp{getline} lee hasta el salto de línea, incluyendo los espacios.
  Regla práctica: si los datos pueden contener espacios, usa \cpp{getline};
  si son tokens separados por espacio, usa \cpp{>>}.

  \medskip
  \textbf{¿Se puede mezclar \cpp{getline} y \cpp{>>} en el mismo stream?}\\
  Sí, pero con cuidado.  \cpp{>>} deja el \texttt{\textbackslash n} en el buffer;
  el siguiente \cpp{getline} lo consumirá como una línea vacía.  Si vas a
  mezclarlos, agrega \cpp{stream.ignore()} o \cpp{stream.ignore(1000,'\textbackslash n')}
  después de cada bloque de \cpp{>>} para limpiar el buffer antes de
  \cpp{getline}.  En la actividad de hoy se evita este problema usando
  \cpp{getline} para cadenas y \cpp{>>} para números en campos separados.

  \medskip
  \textbf{¿Qué es EOF exactamente?}\\
  \textit{End Of File}.  No es un byte especial dentro del archivo; es
  una condición que el sistema operativo reporta cuando el stream llega
  al final de los datos.  En el contexto del teclado, \texttt{Ctrl+D}
  (Linux) o \texttt{Ctrl+Z} (Windows) simula un EOF.

  \medskip
  \textbf{¿\cpp{endl} vs \cpp{"\textbackslash n"} al escribir en archivo?}\\
  \cpp{endl} escribe \texttt{\textbackslash n} \textit{y} vacía el buffer
  (\textit{flush}).  \cpp{"\textbackslash n"} solo escribe el carácter.  Para
  archivos, \cpp{"\textbackslash n"} es más rápido si escribes muchas líneas.
  \cpp{endl} es más seguro si necesitas garantizar que el dato llegó al
  disco inmediatamente (por ejemplo, en un log de errores).

\end{notaprofesor}

% =============================================================================
\section{Actividad: agenda de cumpleaños}
% =============================================================================

\begin{notaprofesor}[Organización de la actividad (35 min)]

  \textbf{Antes de la clase:}
  \begin{itemize}
    \item Crear la carpeta en Google Drive y tener el link listo.
    \item Preparar un archivo \texttt{archivos.txt} de ejemplo (vacío al
          inicio, se llena durante la Parte~A) que los alumnos puedan
          usar para iterar sobre todos los archivos individuales.
  \end{itemize}

  \textbf{Parte~A (10 min) — Cada alumno genera su archivo:}
  \begin{enumerate}
    \item Dar el link de Drive.
    \item Los alumnos escriben un programa con \cpp{ofstream} que genere su
          archivo personal.  El formato es cuatro líneas: nombre, apellidos,
          día, mes.
    \item Suben el archivo generado (no el \texttt{.cpp}) a Drive.
    \item Mientras suben, actualizar \texttt{archivos.txt} en Drive con la
          lista de nombres de archivo para que todos puedan descargarlo.
  \end{enumerate}

  \textbf{Parte~B (25 min) — Las consultas:}\\
  Dar tiempo libre.  Circular por el salón para desbloquear a quienes
  estén atascados.  Los puntos de fricción más comunes son:
  \begin{itemize}
    \item Consulta~1: cómo leer un archivo cuyo nombre viene de otro archivo
          (abrir un \cpp{ifstream} con un \cpp{string} como nombre).
    \item Consulta~3: cómo manejar el ``año circular'' (si ya pasó la fecha
          este año, buscar la más cercana para el año siguiente).
    \item Consulta~4: cómo representar la ``distancia'' entre dos fechas
          para poder compararlas.
  \end{itemize}

  \textbf{Cierre (5 min):}
  Pedir a 2--3 equipos que muestren su solución a una de las consultas.
  Elegir preferentemente la~3 o la~4, que tienen mayor riqueza algorítmica.

\end{notaprofesor}

\begin{notaprofesor}[Solución de referencia — Consulta 1: unir archivos]
\begin{verbatim}
#include <iostream>
#include <fstream>
#include <string>
using namespace std;

int main(){
    ifstream lista("archivos.txt");
    ofstream agenda("agenda.txt");
    string archivo;

    while(getline(lista, archivo)){
        ifstream f(archivo);
        string nombre, apellido;
        int dia, mes;
        getline(f, nombre);
        getline(f, apellido);
        f >> dia >> mes;
        agenda << nombre << " " << apellido
               << " " << dia << " " << mes << endl;
        f.close();
    }
    agenda.close();
}
\end{verbatim}
  Nota: \cpp{ifstream f(archivo)} acepta un \cpp{string} como nombre
  desde C++11.  En compiladores más antiguos puede requerirse
  \cpp{f.open(archivo.c\_str())}.
\end{notaprofesor}

\begin{notaprofesor}[Solución de referencia — Consulta 2: filtrar por mes]
\begin{verbatim}
#include <iostream>
#include <fstream>
#include <string>
using namespace std;

int main(){
    int mes_buscado;
    cout << "Mes (1-12): ";
    cin >> mes_buscado;

    ifstream agenda("agenda.txt");
    string nombre, apellido;
    int dia, mes;

    while(agenda >> nombre >> apellido >> dia >> mes){
        if(mes == mes_buscado)
            cout << nombre << " " << apellido << endl;
    }
}
\end{verbatim}
\end{notaprofesor}

\begin{notaprofesor}[Solución de referencia — Consulta 3: próximo cumpleaños]
  La idea es convertir cada cumpleaños a ``día del año'' y comparar con
  el día actual.  Se puede usar mes$\times$31~+~día como aproximación
  suficiente para encontrar el mínimo; o calcular el día del año real
  acumulando días por mes.

  Estrategia simple con día del año aproximado:
\begin{verbatim}
// dias_aprox: aproximación suficiente para ordenar
int dias_aprox(int dia, int mes){
    int acum[] = {0,31,59,90,120,151,181,212,243,273,304,334};
    return acum[mes-1] + dia;
}

// HOY = 6 de mayo = día ~126
int hoy = dias_aprox(6, 5);

// Para cada persona calcular dias_aprox(dia, mes)
// Si > hoy: candidato de este año
// Si <= hoy: candidato del año siguiente (sumar 365)
// El próximo cumpleaños es el candidato mínimo
\end{verbatim}
  En clase es válido que los alumnos elijan la aproximación mes$\times$31~+~día
  si les funciona para ordenar.  El objetivo es algorítmico, no calendárico.
\end{notaprofesor}

\begin{notaprofesor}[Solución de referencia — Consulta 4: par más cercano]
  Convertir todas las fechas a días del año (misma función que consulta 3),
  ordenarlas de menor a mayor, y buscar la diferencia mínima entre
  elementos consecutivos.  No olvidar la diferencia circular entre el
  último y el primero (suma 365 al último si es menor que el primero).
\begin{verbatim}
// Con N personas y arreglos nombre[], dia_año[]
// Ordenar por dia_año (bubble sort o cualquier método conocido)
// Luego:
int min_diff = 365;
int p1 = 0, p2 = 1;
for(int i = 0; i < N-1; i++){
    int diff = dia_año[i+1] - dia_año[i];
    if(diff < min_diff){
        min_diff = diff;
        p1 = i; p2 = i+1;
    }
}
// Diferencia circular:
int diff_circular = 365 - dia_año[N-1] + dia_año[0];
if(diff_circular < min_diff){
    min_diff = diff_circular;
    p1 = N-1; p2 = 0;
}
\end{verbatim}
\end{notaprofesor}

% =============================================================================
\section{Errores frecuentes}
% =============================================================================

\begin{notaprofesor}[Los cinco errores más comunes]

  \textbf{1.\ Archivo no encontrado — sin mensaje de error.}\\
  Síntoma: el programa termina sin hacer nada.  No hay excepción por defecto.\\
  Remedio: siempre verificar \cpp{if (!ifile)} justo después de abrir.

  \vspace{0.5em}
  \textbf{2.\ Ejecutable en directorio distinto al archivo de datos.}\\
  Síntoma: mismo error que el anterior.\\
  Remedio: usar una ruta relativa correcta o ejecutar el programa desde
  el directorio donde están los datos.  En VS Code el directorio de trabajo
  del ejecutable puede ser distinto al del \texttt{.cpp}.

  \vspace{0.5em}
  \textbf{3.\ Mezclar \cpp{>>} y \cpp{getline} sin limpiar el buffer.}\\
  Síntoma: \cpp{getline} lee una cadena vacía en la siguiente llamada.\\
  Remedio: agregar \cpp{stream.ignore()} o \cpp{stream.ignore(1000,'\textbackslash n')}
  después de cada bloque de \cpp{>>}.

  \vspace{0.5em}
  \textbf{4.\ Olvidar \cpp{close()} en \cpp{ofstream}.}\\
  Síntoma: el archivo existe pero está vacío o incompleto.  El buffer no
  se vació.\\
  Remedio: llamar \cpp{close()} explícitamente, o mejor aún, envolver el
  \cpp{ofstream} en un bloque \cpp{\{\}} para que su destructor lo cierre.

  \vspace{0.5em}
  \textbf{5.\ Abrir un \cpp{ofstream} para leer (o viceversa).}\\
  Síntoma: error de compilación poco claro, o el archivo se borra al abrirlo
  con \cpp{ofstream} (que crea/sobreescribe).\\
  Remedio: \texttt{i}~=~input (leer), \texttt{o}~=~output (escribir).

\end{notaprofesor}

% =============================================================================
\section{Soluciones a los ejercicios}
% =============================================================================

\begin{notaprofesor}[Ejercicio 1 — Calificaciones]
\begin{verbatim}
#include <iostream>
#include <fstream>
using namespace std;

int main(){
    ofstream fout("calificaciones.txt");
    float c1, c2, c3;
    cout << "Tres calificaciones: ";
    cin >> c1 >> c2 >> c3;
    fout << c1 << endl << c2 << endl << c3 << endl;
    fout.close();

    ifstream fin("calificaciones.txt");
    float x, sum = 0;
    int n = 0;
    while(fin >> x){ sum += x; n++; }
    cout << "Promedio: " << sum/n << endl;
}
\end{verbatim}
\end{notaprofesor}

\begin{notaprofesor}[Ejercicio 2 — Agregar calificaciones]
  Solo cambia la línea de apertura de \cpp{ofstream}:
  \begin{verbatim}
ofstream fout("calificaciones.txt", fstream::app);
  \end{verbatim}
  El bloque de lectura y promedio queda exactamente igual que en el
  ejercicio 1.  Al acumular todas las calificaciones del archivo completo
  el promedio ya incluye las anteriores.
\end{notaprofesor}

\begin{notaprofesor}[Ejercicio 3 — Canciones numeradas]
\begin{verbatim}
#include <iostream>
#include <fstream>
#include <string>
using namespace std;

int main(){
    ifstream fin("canciones.txt");
    ofstream fout("canciones_numeradas.txt");
    string linea;
    int num = 1;
    while(getline(fin, linea)){
        fout << num << ". " << linea << endl;
        num++;
    }
}
\end{verbatim}
\end{notaprofesor}

\begin{notaprofesor}[Ejercicio 4 — Concatenar archivos]
\begin{verbatim}
#include <iostream>
#include <fstream>
#include <string>
using namespace std;

int main(){
    string a, b, c;
    cout << "Primer archivo: ";  cin >> a;
    cout << "Segundo archivo: "; cin >> b;
    cout << "Archivo resultado: "; cin >> c;

    ifstream f1(a), f2(b);
    ofstream fout(c);
    string linea;
    while(getline(f1, linea)) fout << linea << endl;
    while(getline(f2, linea)) fout << linea << endl;
}
\end{verbatim}
\end{notaprofesor}

\begin{notaprofesor}[Ejercicio 5 — Archivo inexistente]
  Al abrir un \cpp{ifstream} con un nombre de archivo que no existe:
  \begin{itemize}
    \item No hay excepción por defecto.
    \item No hay error de compilación.
    \item El stream queda en estado de error (\cpp{failbit} activado).
    \item Cualquier operación de lectura sobre él no hace nada.
    \item \cpp{getline} devuelve inmediatamente \cpp{false} → el
          \cpp{while} no itera.
    \item \cpp{>>} no modifica la variable de destino (queda con su
          valor anterior, normalmente basura).
  \end{itemize}
  Para detectar el error: \cpp{if (!fin)} o \cpp{if (fin.fail())} justo
  después de abrir.  Para activar excepciones explícitas:
  \begin{verbatim}
fin.exceptions(ifstream::failbit | ifstream::badbit);
  \end{verbatim}
  Después de eso, abrir un archivo que no existe lanza
  \cpp{std::ios\_base::failure}.
\end{notaprofesor}

% =============================================================================
\section{Rúbrica de la actividad}
% =============================================================================

\begin{center}
\begin{tabular}{p{7.5cm}cc}
  \toprule
  \textbf{Criterio} & \textbf{Puntos} & \textbf{Evidencia} \\
  \midrule
  \multicolumn{3}{l}{\textit{Parte~A — Generar archivo personal}} \\
  \quad Archivo con 4 líneas en el formato correcto  & 5 & archivo en Drive \\
  \quad Generado desde un programa C++ con ofstream  & 5 & código fuente \\
  \midrule
  \multicolumn{3}{l}{\textit{Consulta 1 — Unir archivos}} \\
  \quad Lee correctamente cada archivo individual     & 5 & agenda.txt generada \\
  \quad agenda.txt tiene una persona por línea        & 5 & formato correcto \\
  \midrule
  \multicolumn{3}{l}{\textit{Consulta 2 — Filtrar por mes}} \\
  \quad Lee agenda.txt con ifstream                   & 5 & código fuente \\
  \quad Muestra solo las personas del mes pedido      & 5 & salida correcta \\
  \midrule
  \multicolumn{3}{l}{\textit{Consulta 3 — Próximo cumpleaños}} \\
  \quad Maneja correctamente el caso ``ya pasó este año'' & 5 & prueba con fecha pasada \\
  \quad Muestra todos los compañeros si coinciden      & 5 & prueba con empate \\
  \midrule
  \multicolumn{3}{l}{\textit{Consulta 4 — Par más cercano}} \\
  \quad Resultado correcto para el conjunto completo   & 5 & salida verificada \\
  \quad Maneja el caso circular (dic--ene)             & 5 & prueba con fechas extremas \\
  \midrule
  \textbf{Total}                                       & \textbf{50} & \\
  \bottomrule
\end{tabular}
\end{center}

% =============================================================================
\section{Conexiones con temas posteriores}
% =============================================================================

\begin{itemize}
  \item \textbf{Sesión 22 — Archivos binarios:} el mismo \cpp{<fstream>},
        pero con el flag \cpp{fstream::binary} y los métodos \cpp{write}/\cpp{read}.
        Los alumnos ya entienden el modelo de stream; solo cambia cómo se
        serializan los datos.
  \item \textbf{Persistencia en POO:} en el curso de Orientación a Objetos
        los objetos se guardarán y recuperarán de archivos.  El patrón
        write/read que ven hoy es la base de toda serialización.
  \item \textbf{Bases de datos:} un DBMS como SQLite es, al final, un
        gestor de archivos muy sofisticado.  Las operaciones CRUD
        (Create, Read, Update, Delete) mapean directamente a lo que
        hicieron hoy con \cpp{ofstream} e \cpp{ifstream}.
\end{itemize}

% =============================================================================
\section{Materiales y recursos}
% =============================================================================

\begin{itemize}
  \item \textbf{Notas alumno:}
        \texttt{notas\_alumno/pdf/sesion21\_archivos\_texto.pdf}
  \item \textbf{Material inicial:}
        \texttt{materialInicial/sesion21/} — \texttt{archivos1.cpp},
        \texttt{archivos2.cpp}, \texttt{archivos3.cpp}, \texttt{canciones.txt},
        \texttt{resultados.txt}
  \item \textbf{Hex editor en línea:} \texttt{hexed.it}
  \item \textbf{Google Drive:} carpeta compartida para la actividad
        (crear antes de la sesión y tener el link listo)
\end{itemize}

\end{document}
